% Complete chapter-1 typesetting source; compact textbook presentation.
% Authoritative Markdown: ../../chapters/ch01/student.md
% Stable task/source IDs are comments. All learner prose, tasks, and notes are retained.
% Shared layout for the two complete chapter-1 PDFs.
% This is a small hand-maintained TeX layout, not a conversion pipeline.
\RequirePackage{fix-cm}
\documentclass[11pt,twoside]{article}
\usepackage{fontspec}
\setmainfont{Cambria}[
  Renderer=HarfBuzz,
  BoldFont=cambriab.ttf,
  ItalicFont=cambriai.ttf,
  BoldItalicFont=cambriaz.ttf
]
\setmonofont{lmmono10-regular.otf}[BoldFont=lmmonolt10-bold.otf]
\usepackage[paperwidth=165mm,paperheight=241mm,inner=21mm,outer=19mm,
  top=22mm,bottom=20mm,headheight=14pt,headsep=7mm,footskip=10mm]{geometry}
\usepackage{xcolor,array,booktabs,longtable,tabularx}
\usepackage{enumitem,fancyhdr}
\usepackage[unicode,hidelinks]{hyperref}
\definecolor{ink}{HTML}{000000}
\definecolor{muted}{HTML}{4C5558}
\setlength{\parindent}{0pt}
\setlength{\parskip}{3.5pt}
\AtBeginDocument{\fontsize{11}{14.3}\selectfont}
\setlength{\tabcolsep}{4pt}
\renewcommand{\arraystretch}{1.10}
\setlist[itemize]{leftmargin=1.3em,label=\textbullet,itemsep=2pt,topsep=3pt,parsep=0pt}
\setlist[enumerate]{leftmargin=1.8em,itemsep=3pt,topsep=4pt,parsep=0pt}
\widowpenalty=10000
\clubpenalty=10000
\emergencystretch=1em
\raggedbottom
\pagestyle{fancy}
\fancyhf{}
\fancyhead[LE,RO]{\small\itshape Naming the world}
\fancyhead[LO,RE]{\small Cook Expanded}
\fancyfoot[LE,RO]{\small\thepage}
\renewcommand{\headrulewidth}{0.3pt}
\newcommand{\booktitle}[2]{
  {\fontsize{9}{11}\selectfont\color{muted}#1\par}
  \vspace{3pt}
  {\fontsize{20}{24}\selectfont\color{ink}#2\par}
  \vspace{6pt}
}
\newcommand{\source}[1]{\par{\footnotesize\color{muted}#1\par}}
\newcommand{\readingtext}[1]{
  \par\addvspace{5pt}
  \begingroup\leftskip=1.5em
    \fontsize{12}{16}\selectfont #1\par
  \endgroup\addvspace{5pt}
}
\makeatletter
\renewcommand{\section}{\@startsection{section}{1}{\z@}
  {-2.4ex \@plus -0.6ex \@minus -0.2ex}{0.8ex \@plus 0.2ex}
  {\normalfont\fontsize{12}{15}\selectfont\bfseries\color{ink}}}
\renewcommand{\subsection}{\@startsection{subsection}{2}{\z@}
  {-1.8ex \@plus -0.5ex \@minus -0.2ex}{0.5ex \@plus 0.2ex}
  {\normalfont\fontsize{11}{14}\selectfont\bfseries\color{ink}}}
\makeatother

\hypersetup{pdftitle={Chapter 1: Naming the world},pdfauthor={Cook Expanded}}
\fancyfoot[LO,RE]{\footnotesize\color{muted}Chapter 1}
\begin{document}
\booktitle{CHAPTER 1}{Naming the world}


% Chapter: ch01. Compact presentation revision awaiting review. Source crosswalk: coverage.json.


\pdfbookmark[1]{1. Orthography}{section1}
\section*{1. Orthography}


\begingroup\small
\begin{longtable}{@{}>{\raggedright\arraybackslash}p{35mm} >{\raggedright\arraybackslash}p{87mm}@{}}
\toprule
\textbf{Form} & \textbf{Value or convention} \\
\midrule
\endfirsthead
\multicolumn{2}{@{}l}{\small\itshape Continued} \\
\toprule
\textbf{Form} & \textbf{Value or convention} \\
\midrule
\endhead
\midrule
\multicolumn{2}{r@{}}{\footnotesize Continued on the next page} \\
\endfoot
\bottomrule
\endlastfoot
\textbf{æ} & \textit{Ash}, a vowel letter distinct from \textbf{a} and \textbf{e}. \\
\textbf{ð, þ} & \textit{Eth} and \textit{thorn}, representing the sounds of English \textbf{th}. \\
\textbf{ā, ē, ī, ō, ū, ȳ, ǣ} & A macron marks vowel length, not stress. In diphthongs it stands on the first element: \textbf{trēo}. \\
\textbf{ę} & Hooked \textbf{e} marks historical vowel origin: \textbf{męte}. It is not a length mark or an inflection. \\
\end{longtable}
\endgroup


Retain the spelling and vowel marks given here.


\pdfbookmark[1]{2. Nouns and the article}{section2}
\section*{2. Nouns and the article}


Nouns are masculine, feminine, or neuter; gender must be learned with each noun. The article agrees in gender, number, and case. The nominative is the case of the subject and is used for the isolated phrases in this lesson.


\begingroup\small
\begin{longtable}{@{}>{\raggedright\arraybackslash}p{32mm} >{\raggedright\arraybackslash}p{29mm} >{\raggedright\arraybackslash}p{29mm} >{\raggedright\arraybackslash}p{26.5mm}@{}}
\toprule
\textbf{Number} & \textbf{Masculine} & \textbf{Feminine} & \textbf{Neuter} \\
\midrule
\endfirsthead
\multicolumn{4}{@{}l}{\small\itshape Continued} \\
\toprule
\textbf{Number} & \textbf{Masculine} & \textbf{Feminine} & \textbf{Neuter} \\
\midrule
\endhead
\midrule
\multicolumn{4}{r@{}}{\footnotesize Continued on the next page} \\
\endfoot
\bottomrule
\endlastfoot
Singular & \textbf{se} & \textbf{sēo} & \textbf{ðæt} \\
Plural & \textbf{ðā} & \textbf{ðā} & \textbf{ðā} \\
\end{longtable}
\endgroup


These nominative forms of the demonstrative serve here as the definite article, “the.”


% example: ch01-example-01


\textbf{se dæg}, the day; \textbf{sēo niht}, the night; \textbf{ðæt sǣd}, the seed.


% example: ch01-example-03


\textbf{ðā ðīestru}, the darkness: \textbf{ðīestru} is plural, despite its singular English equivalent. It is the only plural noun form introduced here; no general rule of plural formation is implied.


\pdfbookmark[1]{3. Coordination}{section3}
\section*{3. Coordination}


\textbf{and}, “and,” is invariable. Each article agrees with its own noun; coordination does not change singular articles to plural.


% example: ch01-example-02


\textbf{se fisc and sēo wyrt}, the fish and the herb.


\pdfbookmark[1]{4. Vocabulary}{section4}
\section*{4. Vocabulary}


Nominative singular unless marked plural. \textbf{m.} masculine; \textbf{f.} feminine; \textbf{n.} neuter. Other case forms are not introduced here.


\begingroup\small
\begin{longtable}{@{}>{\raggedright\arraybackslash}p{30mm} >{\raggedright\arraybackslash}p{63mm} >{\raggedright\arraybackslash}p{26mm}@{}}
\toprule
\textbf{Word} & \textbf{Meaning} & \textbf{Gender / number} \\
\midrule
\endfirsthead
\multicolumn{3}{@{}l}{\small\itshape Continued} \\
\toprule
\textbf{Word} & \textbf{Meaning} & \textbf{Gender / number} \\
\midrule
\endhead
\midrule
\multicolumn{3}{r@{}}{\footnotesize Continued on the next page} \\
\endfoot
\bottomrule
\endlastfoot
\textbf{ǣfen} & evening & n. \\
\textbf{morgen} & morning & m. \\
\textbf{dæg} & day & m. \\
\textbf{niht} & night & f. \\
\textbf{eorðe} & earth & f. \\
\textbf{heofon} & heaven; sky & m. \\
\textbf{fisc} & fish & m. \\
\textbf{fugol} & bird & m. \\
\textbf{gærs} & grass & n. \\
\textbf{sǣd} & seed & n. \\
\textbf{trēo} & tree & n. \\
\textbf{wyrt} & herb; plant & f. \\
\textbf{drȳgnes} & dry land & f. \\
\textbf{wæter} & water & n. \\
\textbf{andlīcnes} & image & f. \\
\textbf{gelīcnes} & likeness & f. \\
\textbf{lēoht} & light & n. \\
\textbf{ðīestru} & darkness & plural noun \\
\textbf{tācen} & sign & n. \\
\textbf{tīd} & season; time & f. \\
\textbf{gāst} & spirit & m. \\
\textbf{līf} & life & n. \\
\textbf{gesceaft} & creature; created being or thing & f. \\
\textbf{męte} & food & m. \\
\textbf{gēar} & year & n. \\
\textbf{mann} & man; human being & m. \\
\textbf{dēor} & animal; beast & n. \\
\textbf{wæstm} & fruit; produce & m. \\
\textbf{fæstnes} & firmament; heavenly expanse & f. \\
\textbf{nēowolnes} & abyss; the deep & f. \\
\end{longtable}
\endgroup


\clearpage
\pdfbookmark[1]{5. Preparation}{section5}
\section*{5. Preparation}


\begin{enumerate}
% item: ch01-prep-01
\item[1.] Supply the nominative singular article: \textbf{\_\_\_ eorðe}.
% item: ch01-prep-02
\item[2.] Write “the evening” with the gender selected in §4; state the gender.
% item: ch01-prep-03
\item[3.] Supply the article and identify the number: \textbf{\_\_\_ ðīestru}.
\end{enumerate}


\pdfbookmark[1]{6. Composition}{section6}
\section*{6. Composition}


Translate into Old English. Use the vocabulary of §4, nominative forms, an article before each noun, and \textbf{and}. All nouns are singular except “darkness” in 9.


\begin{enumerate}
% item: EX-I-01-01
\item[1.] The evening and the morning.
% item: EX-I-01-02
\item[2.] The day and the night.
% item: EX-I-01-03
\item[3.] The earth and the heaven.
% item: EX-I-01-04
\item[4.] The fish and the bird.
% item: EX-I-01-05
\item[5.] The grass and the seed.
% item: EX-I-01-06
\item[6.] The tree and the herb.
% item: EX-I-01-07
\item[7.] The dry land and the water.
% item: EX-I-01-08
\item[8.] The image and the likeness.
% item: EX-I-01-09
\item[9.] The light and the darkness (plur.).
% item: EX-I-01-10
\item[10.] The sign and the season.
% item: EX-I-01-11
\item[11.] The spirit and the life.
% item: EX-I-01-12
\item[12.] The creature and the food.
% item: EX-I-01-13
\item[13.] The year and the day.
% item: EX-I-01-14
\item[14.] The man and the animal.
% item: EX-I-01-15
\item[15.] The tree and the fruit.
% item: EX-I-01-16
\item[16.] The firmament and the abyss.
\end{enumerate}


\clearpage
\pdfbookmark[1]{7. Reading and translation}{section7}
\section*{7. Reading and translation}


A clause from an account of the Creation. The semicolon marks the continuation of the original sentence.


The following forms are supplied for these tasks, not assigned for independent recall.


\begingroup\small
\begin{longtable}{@{}>{\raggedright\arraybackslash}p{23mm} >{\raggedright\arraybackslash}p{99mm}@{}}
\toprule
\textbf{Form} & \textbf{Meaning and analysis} \\
\midrule
\endfirsthead
\multicolumn{2}{@{}l}{\small\itshape Continued} \\
\toprule
\textbf{Form} & \textbf{Meaning and analysis} \\
\midrule
\endhead
\midrule
\multicolumn{2}{r@{}}{\footnotesize Continued on the next page} \\
\endfoot
\bottomrule
\endlastfoot
\textbf{sōðlice} & now; for its part. Connective, not a present-time expression here. \\
\textbf{wæs} & was. Third-person singular past indicative of \textbf{wesan}, “be”; linking verb. \\
\textbf{īdel} & empty; void. Predicate adjective, feminine nominative singular, describing \textbf{eorðe}. \\
\textbf{ǣmtig} & empty; void. Predicate adjective, feminine nominative singular, also describing \textbf{eorðe}. \\
\end{longtable}
\endgroup


% item: ch01-text-01


\readingtext{Sēo eorðe sōðlice wæs īdel and ǣmtig;}


Using the supplied forms, answer in English:


\begin{enumerate}
% item: ch01-reading-01
\item[1.] What is described, and in what state?
% item: ch01-reading-02
\item[2.] Identify the subject's article and noun; give gender, number, and case. What does \textbf{and} join?
% item: ch01-reading-03
\item[3.] Translate the clause.
\end{enumerate}


\pdfbookmark[1]{8. Review}{section8}
\section*{8. Review}


For 1–2, use the taught vocabulary, nominative singular forms, an article before each noun, and \textbf{and}.


\begin{enumerate}
% item: ch01-review-01
\item[1.] The seed and the herb.
% item: ch01-review-02
\item[2.] The heaven and the water.
% item: ch01-review-03
\item[3.] Explain the articles in \textbf{ðā ðīestru} and \textbf{ðæt lēoht}. Does \textbf{ðīestru} establish a general rule for forming noun plurals?
\end{enumerate}


\begingroup\small
\clearpage
\pdfbookmark[1]{Notes on sources and editorial choices}{section9}
\section*{Notes on sources and editorial choices}


\textbf{Ex:} A. S. Cook, \textit{Exercises in Old English} (1895). \textbf{FB:} \textit{A First Book in Old English}, 2nd ed. (1895). Page references are to printed pages; scan locators and full evidence are in the editorial record.


\begin{enumerate}
\item[1.] \textbf{Orthography (§1):} FB §§2–4, pp.5–6. Source letters and vowel marks are retained, including historical \textbf{ę}.
\item[2.] \textbf{Grammar (§§2–3):} FB §§38–40, pp.26–27; §84, p.50; §149, p.88. The examples are new teaching material.
\item[3.] \textbf{Vocabulary (§4):} Ex, English–Old English vocabulary, corresponding entries; FB §§43–44, 46–47, 51–53. The key records valid alternatives to the selected forms and genders. \textbf{andlīcnes} selects the \textbf{d} permitted by Ex \textbf{an(d)līcnes}.
\item[4.] \textbf{Correction C01 (§4; composition 15):} \textbf{wæstm} replaces Cook's \textbf{wǣstm}. Short /æ/ is supported by \href{https://glossary.oldenglishaerobics.net/oeaginfo.php?objid=xwaestm\&objtype=gentry}{Baker, \textit{wæstm}}, and \href{https://doi.org/10.17811/selim.29.2024.109-25}{Stenroos (2024), p.118, Table 1}. The original is retained in the editorial record.
\item[5.] \textbf{Exercises (§§5–8):} The sixteen composition prompts preserve Ex, Group I, Exercise 1, p.1, including \textbf{(plur.)} in 9. Instructions, additional tasks, examples, and model answers are new teaching material.
\item[6.] \textbf{Reading (§7):} FB, Reader I, p.123, attributed to Ælfric's Genesis: the second sentence's opening clause. Wording, marks, and semicolon are retained; note cues 4–6 and a line break are omitted. Note 5 glosses \textbf{sōðlice} with Latin \textit{autem}. Supplied forms: FB §138, pp.78–79; §§149–150, p.88; §57.3, p.39; vocabulary pp.274, 301, 313. The translation is independently prepared.
\end{enumerate}



\endgroup
\end{document}
